\documentclass{article}

\usepackage[utf8]{inputenc}     % pdfLaTeX
\usepackage[T1]{fontenc}

\usepackage{amsmath}
\usepackage{amsthm}
\usepackage{amssymb}

\title{Applied Algebra}
\author{Rovenszky Robin}
\date{Last updated: 12 March 2026}

\newtheorem{theorem}{Theorem}

\begin{document}

\maketitle

\section{The Binomial Theorem}
\begin{theorem}[Binomial Theorem]
The Binomial theorem describes the algebraic expansion of powers of a binomial.
\begin{equation}
    (x+y)^n = \binom{n}{0}x^ny^0 + \binom{n}{1}x^{n-1}y^1 + ... + \binom{n}{n}x^0y^n.
\end{equation}
\end{theorem}

\section{Problem set.}
\subsection{Algebra 101.}
\subsubsection{Problem 1.1.}
Expand $(a\pm b)^5$.
\begin{equation}
    (a \pm b)^5
\end{equation}
Using the Binomial Theorem (Theorem 1) yields:
\begin{equation}
    (a \pm b)^5 = \binom{5}{0}a^5 \pm \binom{5}{1}a^4b + \binom{5}{2}a^3b^2 \pm \binom{5}{3}a^2b^3 + \binom{5}{4}ab^4 \pm \binom{5}{5}b^5.
\end{equation}

\subsubsection{Problem 1.2.}
Expand $(x\pm 1)^5$.
\begin{equation}
    (x \pm 1)^5
\end{equation}
Using the Binomial Theorem (Theorem 1) yields:
\begin{equation}
    (x \pm 1)^5 = \binom{5}{0}x^5\pm \binom{5}{1}x^4 + \binom{5}{2}x^3 \pm \binom{5}{3}x^2 + \binom{5}{4}x \pm \binom{5}{5}.
\end{equation}

\subsubsection{Problem 1.3.}
Expand $(a+b+c+d)^2$.
\begin{equation}
    (a+b+c+d)^2 = a^2 + b^2 + c^2 + d^2 + 2ab + 2ac + 2ad + 2bc + 2bd + 2cd.
\end{equation}

\subsubsection{Problem 1.4.}
Rewrite $a^2 + ab + a + b$ as a product of two first-degree polynomials.
\begin{equation}
    a^2 + ab + a + b = a(a+b) + (a+b) = (a + b) (a + 1).
\end{equation}

\subsubsection{Problem 1.5.}
Rewrite $x^3 + 2x^2 + 2x + 1$ as a product of a second and a first degree polynomial. (It cannot be rewritten as a product of three first-degree polynomials, due to The Fundamental Theorem of Algebra and results from Abstract Algebra / Field Theory).

\begin{equation}
    x^3 + 2x^2 + 2x + 1 = (x^3 + x^2 + x) + (x^2 + x + 1) = x(x^2 + x + 1) + (x^2 + x + 1) =
\end{equation}
\begin{equation}
    = (x^2 + x + 1)(x + 1).
\end{equation}

\subsubsection{Problem 1.6.}
Rewrite $x^3 + 6x^2 + 12x + 8$ as a product of polynomials.
\\
Method 1.\\
Notice that the coefficients 1, 6, 12, 8 match the binomial expansion pattern of a perfect cube. Then use the identity
\begin{equation}
    a^3 + 3a^2b+3ab^2+b^3 = (a+b)^3.
\end{equation}
Thus,
\begin{equation}
    x^3 + 6x^2 + 12x + 8 = (x+2)^3.
\end{equation}

Method 2.\\
We observe that $-2$ is a root of our polynomial (i.e. a point at which the expression evaluates to $0$), and thus $(x+2)$ is a factor. Dividing the polynomial by $(x+2)$ gives
\begin{equation}
    x^2 + 4x + 4 = (x+2)^2,
\end{equation}
Therefore
\begin{equation}
    x^3 + 6x^2 + 12x + 8 = (x+2)(x+2)^2 = (x+2)^3.
\end{equation}

\subsubsection{Problem 1.7.}
Find all solutions of the equation $ab + 2a - 3b = 14$ where $a, b \in \mathbb{N^+}$.

Group terms containing $a$
\[
    ab + 2a - 3b = 14
\]
\[
    a(b+2) - 3b = 14
\]
Move the $3b$ term
\[
    a(b+2) = 3b + 14
\]
Express $a$:
\[
    a = \frac{3b+14}{b+2} 
\]
Rewrite the numerator:
\[
    3b+14 = 3(b+2) + 8
\]
Thus
\[
    a = 3 + \frac{8}{b+2}.
\]

For $a$ to be integer, $b+2$ must divide 8.\\
Divisors of 8: $1, 2, 4, 8$. Since $b \ge 1$, $b+2 \ge 3$, we are left with two possibilities:
\begin{itemize}
    \item $b + 2 = 8 \implies b = 2$
    \item $b + 2 = 8 \implies b = 6$
\end{itemize}

Now we plug these $b$ values back into our expression for $a$ to find the two solutions to our equation: $(5,2)$ and $(4, 6)$.

\subsubsection{Problem 1.8.}
What is the minimal point of the expression $2x^2 + 2y^2 - 2xy + 4x -6y + 8$?

\subsubsection{Problem 1.9.}
Prove the following inequality:
\[
(a+b+c)(\frac{1}{a} + \frac{1}{b} + \frac{1}{c}) \ge 9
\]
where $a, b, c > 0$.

\subsubsection{Problem 2.1.}
Expand $(2x-3)^3$.

\subsubsection{Problem 2.2.}
Expand $(a-1)(b-1)(c-1)$.

\subsubsection{Problem 2.3.}
Rewrite the following expression as a product:
\begin{equation}
    (a+b)^2 - (a-b)^2.
\end{equation}

\subsubsection{Problem 2.4.}
Rewrite the following expression as a product:
\begin{equation}
    x^2y+xy^2-x^2+y^2.
\end{equation}

\subsubsection{Problem 2.5.}
Rewrite $a^4 - 1$ as a product.

Using difference of squares yields:

\subsubsection{Problem 2.6.}
Find all solutions of the equation
\begin{equation}
    4ab-6a+8b = 14.
\end{equation}

\subsubsection{Problem 2.7.}
Find the minimal point of the following expression
\[
5x^2 + 2y^2 -2x +4y + 4xy + 1.
\]

\subsubsection{Problem 2.8.}
Prove the following inequality:
\[
4x+\frac{1}{x} \ge 4
\]
where $x > 0$.

\end{document}